\section{\tr{Fundamentação teórica}{Background}}
\label{sec:fundamentacao}

\subsection{\tr{Modelos de contagem com exposição}{Count models with exposure}}

\pt{%
Seja $Y_{it}$ o número de acidentes do veículo $i$ no mês $t$ e $E_{it}$ sua exposição (horas de operação ou distância). O modelo de Poisson com \textit{offset} supõe $Y_{it} \sim \text{Poisson}(\mu_{it})$, com $\log \mu_{it} = \log E_{it} + \mathbf{x}_{it}^\top \bm{\beta}$, de modo que $\exp(\mathbf{x}_{it}^\top \bm{\beta})$ é uma taxa por unidade de exposição \parencite{cameron2013}. Dados de acidentes costumam ter variância maior que a média (superdispersão), o que leva à binomial negativa, e excesso de zeros, que pode pedir modelos com inflação de zeros \parencite{lambert1992}. A escolha entre essas alternativas em dados de acidentes é discutida em detalhe por \textcite{lord2010}; em segurança viária, a inflação de zeros é criticada por supor um estado de risco nulo sem interpretação física, quando o excesso de zeros costuma vir de exposição pequena ou de heterogeneidade não modelada \parencite{lord2005}.
}
\en{%
Let $Y_{it}$ be the number of accidents of vehicle $i$ in month $t$ and $E_{it}$ its exposure (operating hours or distance). The Poisson model with an offset assumes $Y_{it} \sim \text{Poisson}(\mu_{it})$, with $\log \mu_{it} = \log E_{it} + \mathbf{x}_{it}^\top \bm{\beta}$, so that $\exp(\mathbf{x}_{it}^\top \bm{\beta})$ is a rate per unit of exposure \parencite{cameron2013}. Accident data usually show variance larger than the mean (overdispersion), which leads to the negative binomial, and excess zeros, which may call for zero-inflated models \parencite{lambert1992}. The choice among these alternatives for crash data is discussed in detail by \textcite{lord2010}; in road safety, zero inflation is criticized for assuming a zero-risk state with no physical interpretation, when excess zeros usually come from small exposure or unmodeled heterogeneity \parencite{lord2005}.
}

\subsection{\tr{Estrutura hierárquica e encolhimento}{Hierarchical structure and shrinkage}}

\pt{%
Veículos estão agrupados em unidades, e observações do mesmo veículo se repetem no tempo. Efeitos aleatórios de veículo e de unidade capturam essa dependência e produzem estimativas que encolhem as taxas de grupos pequenos em direção à média geral, na medida da informação disponível \parencite{gelman2007}. É a mesma ideia do Bayes empírico usado em mapeamento de doenças e em segurança viária para identificar locais de alto risco sem ser enganado pela regressão à média \parencite{clayton1987,hauer1997}. Rankings de instituições têm limites conhecidos, porque a incerteza sobre a posição de cada unidade costuma ser grande \parencite{goldstein1996}, e gráficos de funil oferecem uma forma visual de comparar unidades com tamanhos diferentes \parencite{spiegelhalter2005}. Os modelos podem ser ajustados por máxima verossimilhança com aproximação de Laplace \parencite{brooks2017} ou de forma bayesiana \parencite{burkner2017}.
}
\en{%
Vehicles are grouped in sites, and observations of the same vehicle repeat over time. Vehicle and site random effects capture this dependence and yield estimates that shrink the rates of small groups toward the overall mean, to the extent of the available information \parencite{gelman2007}. This is the same idea as the empirical Bayes approach used in disease mapping and in road safety to identify high-risk locations without being misled by regression to the mean \parencite{clayton1987,hauer1997}. Institutional rankings have known limits, because the uncertainty about each unit's position is usually large \parencite{goldstein1996}, and funnel plots offer a visual way to compare units of different sizes \parencite{spiegelhalter2005}. The models can be fitted by maximum likelihood with the Laplace approximation \parencite{brooks2017} or in a Bayesian way \parencite{burkner2017}.
}

\subsection{\tr{Telemetria e risco}{Telematics and risk}}

\pt{%
Sistemas de telemetria registram eventos cinemáticos (acelerações, frenagens e curvas acima de limiares) e eventos de comportamento (velocidade acima do limite, uso de cinto, uso de celular). Estudos de condução naturalística associam a frequência desses eventos ao risco de colisão de cada motorista \parencite{guo2013}. Em caminhões, eventos críticos registrados pela telemetria de uma grande transportadora dos EUA se associaram a acidentes num modelo binomial negativo com a distância como exposição \parencite{cai2021}, e frenagens bruscas por hora previram sinistros com dano material \parencite{driessen2024}. Em pátio, onde as velocidades são baixas e as manobras dominam, não é óbvio quais eventos carregam sinal; o projeto testa isso diretamente (H2).
}
\en{%
Telematics systems record kinematic events (accelerations, braking and turns above thresholds) and behavioral events (speeding, seat belt use, phone use). Naturalistic driving studies associate the frequency of these events with each driver's crash risk \parencite{guo2013}. For trucks, safety-critical events recorded by the telematics of a large US carrier were associated with crashes in a negative binomial model with distance as exposure \parencite{cai2021}, and harsh braking per hour predicted property damage incidents \parencite{driessen2024}. In a yard, where speeds are low and maneuvering dominates, it is not obvious which events carry signal; the project tests this directly (H2).
}

\subsection{\tr{Subnotificação}{Underreporting}}

\pt{%
A subnotificação de lesões de trabalho é grande e varia entre locais: estudos nos EUA que cruzam o levantamento oficial com outros registros encontraram grande parte das lesões fora das estatísticas, com muita variação entre estados \parencite{leigh2004,rosenman2006,boden2008}. Quando o mesmo tipo de evento é registrado por duas fontes imperfeitas, a sobreposição entre elas informa quantos eventos nenhuma das duas captou. É o princípio da captura e recaptura com múltiplas fontes, usado em epidemiologia \parencite{iwgdmf1995}. Aqui, as fontes são os relatos formais e as ordens de manutenção abertas por dano. A hipótese de independência entre as fontes raramente vale na prática: com dependência positiva, como a de acidentes graves que tendem a ser ao mesmo tempo relatados e reparados, o total é subestimado \parencite{hook1995}, e o projeto tratará isso de forma explícita (Seção~\ref{sec:subnotificacao}). A ligação entre registros das duas fontes segue a lógica da ligação probabilística de registros \parencite{fellegi1969}.
}
\en{%
Underreporting of work injuries is large and varies across places: US studies cross-matching the official survey with other records found a large share of injuries missing from the statistics, with much variation across states \parencite{leigh2004,rosenman2006,boden2008}. When the same type of event is recorded by two imperfect sources, the overlap between them informs how many events neither source captured. This is the principle of multiple-source capture-recapture used in epidemiology \parencite{iwgdmf1995}. Here the sources are formal reports and maintenance orders opened for damage. The independence assumption between sources rarely holds in practice: with positive dependence, such as severe accidents tending to be both reported and repaired, the total is underestimated \parencite{hook1995}, and the project will handle this explicitly (Section~\ref{sec:subnotificacao}). Linking records across the two sources follows the logic of probabilistic record linkage \parencite{fellegi1969}.
}

\subsection{\tr{Previsão probabilística}{Probabilistic forecasting}}

\pt{%
Para a extensão opcional, a exposição futura será prevista como distribuição, não como valor pontual, com métodos de séries temporais \parencite{hyndman2021} e avaliada por regras de pontuação próprias, como o CRPS e a perda quantílica \parencite{gneiting2007}.
}
\en{%
For the optional extension, future exposure will be forecast as a distribution rather than a point value, with time series methods \parencite{hyndman2021}, and evaluated with proper scoring rules such as the CRPS and the quantile loss \parencite{gneiting2007}.
}
